\section{Determinants of larger matrices}

The formulas for \(2\times2\) and \(3\times3\) matrices are explicit, but they do not provide a practical definition for every size. We now extend the determinant recursively: a determinant of size \(n\times n\) will be expressed in terms of determinants of size \((n-1)\times(n-1)\).

\subsection{Minors}

Let \(A=(a_{ij})\) be an \(n\times n\) matrix. Remove row \(i\) and column \(j\). The remaining matrix has size \((n-1)\times(n-1)\).

\begin{definitionbox}
The determinant of the matrix obtained by deleting row \(i\) and column \(j\) is called the minor of the entry \(a_{ij}\). It is denoted by
\[
M_{ij}.
\]
\end{definitionbox}

\begin{examplebox}
Let
\[
A=\begin{pmatrix}
2&1&3\\
0&-1&4\\
5&2&1
\end{pmatrix}.
\]
To find \(M_{12}\), delete row \(1\) and column \(2\):
\[
M_{12}
=
\det\begin{pmatrix}0&4\\5&1\end{pmatrix}
=0\cdot1-4\cdot5=-20.
\]
\end{examplebox}

\subsection{Cofactors}

The signs used in the recursive formula alternate in a checkerboard pattern:
\[
\begin{pmatrix}
+&-&+&-&\cdots\\
-&+&-&+&\cdots\\
+&-&+&-&\cdots\\
-&+&-&+&\cdots\\
\vdots&\vdots&\vdots&\vdots&\ddots
\end{pmatrix}.
\]
The sign in position \((i,j)\) is \((-1)^{i+j}\).

\begin{definitionbox}
The cofactor of the entry \(a_{ij}\) is
\[
C_{ij}=(-1)^{i+j}M_{ij}.
\]
\end{definitionbox}

For the previous example,
\[
M_{12}=-20,
\qquad
C_{12}=(-1)^{1+2}M_{12}=20.
\]

\subsection{The recursive definition}

We already defined the determinant of a \(1\times1\) matrix by
\[
\det\begin{pmatrix}a\end{pmatrix}=a.
\]
For \(n\ge2\), we define the determinant by expanding along the first row.

\begin{definitionbox}
Let \(A=(a_{ij})\) be an \(n\times n\) matrix. Its determinant is
\[
\det A
=a_{11}C_{11}+a_{12}C_{12}+\cdots+a_{1n}C_{1n}.
\]
Equivalently,
\[
\det A=\sum_{j=1}^{n}a_{1j}C_{1j}.
\]
\end{definitionbox}

This definition reduces the size by one at every step. Eventually every computation reaches determinants of size \(1\times1\).

For a \(2\times2\) matrix,
\[
A=\begin{pmatrix}a&b\\c&d\end{pmatrix},
\]
the recursive definition gives
\[
\det A
=a\det\begin{pmatrix}d\end{pmatrix}
-b\det\begin{pmatrix}c\end{pmatrix}
=ad-bc.
\]
Thus it agrees with the earlier definition.

For a \(3\times3\) matrix,
\[
A=\begin{pmatrix}
a&b&c\\
d&e&f\\
g&h&i
\end{pmatrix},
\]
expansion along the first row gives
\[
\begin{aligned}
\det A
&=a\begin{vmatrix}e&f\\h&i\end{vmatrix}
-b\begin{vmatrix}d&f\\g&i\end{vmatrix}
+c\begin{vmatrix}d&e\\g&h\end{vmatrix}\\
&=a(ei-fh)-b(di-fg)+c(dh-eg)\\
&=aei+bfg+cdh-ceg-bdi-afh.
\end{aligned}
\]
Thus the recursive definition also agrees with the earlier \(3\times3\) formula.

\subsection{Expansion along any row or column}

Although the determinant was defined using the first row, it can be proved that the same value is obtained by expanding along any row or any column.

\begin{theorembox}
Expansion along row \(i\) gives
\[
\det A=\sum_{j=1}^{n}a_{ij}C_{ij}.
\]
Expansion along column \(j\) gives
\[
\det A=\sum_{i=1}^{n}a_{ij}C_{ij}.
\]
\end{theorembox}

This freedom is computationally important. A row or column containing many zeros usually gives the shortest expansion.

\begin{examplebox}
Consider
\[
B=\begin{pmatrix}
4&0&1&2\\
0&3&0&-1\\
2&0&5&0\\
1&0&0&6
\end{pmatrix}.
\]
The second column contains only one non-zero entry. Expanding along that column,
\[
\det B=3C_{22}.
\]
Since the sign in position \((2,2)\) is positive,
\[
\det B
=3\det\begin{pmatrix}
4&1&2\\
2&5&0\\
1&0&6
\end{pmatrix}.
\]
Expanding the remaining determinant along its second row or using the \(3\times3\) formula gives
\[
\det\begin{pmatrix}
4&1&2\\
2&5&0\\
1&0&6
\end{pmatrix}=98.
\]
Therefore
\[
\det B=294.
\]
\end{examplebox}

\subsection{Using Laplace expansion intelligently}

Laplace expansion is a definition and a general method, but it should not be applied blindly. Before expanding:
\begin{itemize}
\item inspect the matrix for zeros;
\item choose a row or column that produces few non-zero terms;
\item write the cofactor signs carefully;
\item keep the smaller determinants visible;
\item use row or column properties when they shorten the calculation.
\end{itemize}
